\section{Fazit}

Pablo ist ein kompakter Desk-Companion-Bot, der zur Unterstützung von Nutzenden im Alltag entwickelt wurde. 
Das Produkt wurde als physisch präsentes Assistenzsystem konzipiert, das im direkten Umfeld der Nutzenden eingesetzt werden kann und eine niedrigschwellige Interaktion ermöglicht. 
Im Unterschied zu rein softwarebasierten Assistenzlösungen basiert Pablo auf dem Ansatz eines physischen Systems, das sprachliche, visuelle und funktionale Elemente in einem Produkt vereint. 
Ziel des entwickelten Systems ist es, spontane Aufgaben, Fragen und Erinnerungen entgegenzunehmen und in einer möglichst natürlichen Form zu verarbeiten.

Die Interaktion mit dem System erfolgt primär sprachbasiert. 
Dadurch kann die Nutzung ohne klassische Eingabemedien wie Tastatur, Maus oder Smartphone erfolgen. 
Ergänzt wird die sprachliche Interaktion durch visuelle Rückmeldungen über ein Display.
Durch die Kombination aus physischer Präsenz, visueller Darstellung und sprachlicher Kommunikation wurde ein System entwickelt, das nicht ausschließlich funktionale Assistenzaufgaben erfüllt, sondern zugleich als Companion im persönlichen Umfeld wahrgenommen werden kann.

Ein wesentliches Merkmal von Pablo ist der modulare und personalisierbare Systemaufbau. 
Die entwickelte Architektur erlaubt es, einzelne Funktionen bedarfsgerecht zu ergänzen. 
Auf diese Weise kann das Produkt flexibel an unterschiedliche Anforderungen und Nutzungsszenarien angepasst werden. 
Die Personalisierbarkeit beschränkt sich damit nicht nur auf den Funktionsumfang, sondern betrifft grundsätzlich auch die Ausrichtung des Systems auf verschiedene Anwendungskontexte. 
Dadurch ist Pablo nicht ausschließlich für ein einzelnes Einsatzszenario ausgelegt, sondern kann perspektivisch für unterschiedliche Zielgruppen eingesetzt werden, beispielsweise als organisatorische Alltagshilfe, als Lernunterstützung oder als individuell konfigurierbarer Companion im privaten Umfeld. 
Diese Offenheit und Erweiterbarkeit stellt ein zentrales Unterscheidungsmerkmal gegenüber geschlossenen Sprachassistenzsystemen dar, deren Funktionsumfang und Erscheinungsbild in der Regel nur eingeschränkt angepasst werden können.

\section{Ausblick}

Die Entwicklung von Pablo bildet eine Grundlage für weitere technische und konzeptionelle Ausbaustufen. 
Auf Basis des realisierten Prototyps ergeben sich insbesondere Potenziale in den Bereichen Systemstabilität, Funktionsumfang, Personalisierung und Evaluation.

Ein nächster Entwicklungsschritt besteht in der weiteren Optimierung des Gesamtsystems. 
Dazu zählen die Verbesserung des automatischen Systemstarts, die robustere Integration der Hardwarekomponenten sowie die Reduktion von Antwortlatenzen. 
Auch die Abstimmung zwischen Sprachausgabe, visueller Darstellung und mechanischer Bewegung kann weiter verfeinert werden.

Darüber hinaus bietet die modulare Softwarearchitektur eine geeignete Grundlage für die Erweiterung des Funktionsumfangs. 
Künftige Module könnten beispielsweise Erinnerungen, Musiksteuerung, Smart-Home-Anbindungen oder den Zugriff auf externe Wissensquellen umfassen. 
Ebenso erscheint eine weiterentwickelte Plugin-Verwaltung sinnvoll, mit der Funktionen gezielt an unterschiedliche Nutzungsprofile angepasst werden können.

Ein weiteres Entwicklungspotenzial liegt in der stärkeren Personalisierung des Systems, etwa durch adaptive Dialogstrategien, anpassbare Stimmen und visuelle Darstellungen oder lernende Mechanismen zur Berücksichtigung wiederkehrender Routinen.

Neben der technischen Weiterentwicklung ist auch eine systematische Evaluation mit Nutzenden ein wesentlicher nächster Schritt. 
Dabei könnten insbesondere Gebrauchstauglichkeit, Natürlichkeit, Akzeptanz und Vertrauenswürdigkeit des Systems untersucht werden.

Langfristig eröffnet Pablo damit die Perspektive auf ein offenes, anpassbares und physisch verkörpertes Assistenzsystem, das über klassische Sprachassistenz hinausgeht. 
Der entwickelte Prototyp zeigt bereits die grundsätzliche Umsetzbarkeit dieses Ansatzes und bildet eine Basis für eine Weiterentwicklung zu einem vielseitig einsetzbaren Companion-System.
